\documentstyle[%


righttag%


,11pt%


,amssymb%


,russian%               remove in Eng.


,izvru%                 Set izven% in Eng.


]{amsart}





%





\newcommand{\la}{\langle}


\newcommand{\ra}{\rangle}


\newcommand{\Div}{\operatorname{div}}


\newcommand{\Rot}{\operatorname{rot}}


\newcommand{\Grad}{\operatorname{grad}}


\newcommand{\tts}[1]{}


\numberwithin{equation}{section}





\theoremstyle{plain}


\theorembodyfont{\it}


\newtheorem{thms}{\hskip\parindent Теорема}[section]


\newtheorem{lems}{\hskip\parindent Лемма}[section]


%\renewcommand{\thethmnonum}{}





\theoremstyle{definition}


\theorembodyfont{\rm}


\newtheorem{notenonum}{\hskip\parindent Замечание}


\renewcommand{\thenotenonum}{}





%\hoffset=-15mm%                  For Eng.


\setcounter{page}{26}


\god{2004}


\nomer{3~(502)} %                        Remove in Eng.


%\nomer{Vol.\,48, No.\,3}%               Set in Eng.


%\str{\pageref{begin}--\pageref{end}}%  Set in Eng.


\udk{517.983}





\author{\it А.В.\,Калинин, А.А.\,Калинкина}


% NB:  ^^ remove in Eng.


\title{$L_p$-оценки векторных полей}





\date{}


%\pagestyle{myheadings}%         Set in Eng.


\pagestyle{plain} %              Remove in Eng.


\begin{document}


%\thispagestyle{plain}%          Set in Eng.


\maketitle


\label{begin}








При решении математических проблем


гидродинамики и электромагнитной


теории важную роль играют различные оценки, связывающие нормы вектор-функции,


ее ротора и дивергенции в пространствах Лебега. В частности, глубокие


результаты в этом направлении и их приложения к решению конкретных


задач содержатся в 


\cite{1}--\cite{9}. В данной работе исследуется возможность получения


подобных оценок с использованием  некоторых представлений векторных полей,


рассмотренных в \cite{10}.


\addtocounter{section}{1}


\medskip





{\bf 1. Обозначения.}


Пусть $\Omega\in R^3$ --- открытая односвязная


область с регулярной границей $\Gamma$, $\overline{\Omega}$ --- замыкание


$\Omega$ в $R^3$, $\bold n(\bold x)$, $\bold x\in\Gamma$, --- единичный вектор


внешней нормали.


Через ${\cal D}(\Omega)=C_0^{\infty}(\Omega)$ обозначено пространство


пробных функций


${\varphi:\Omega\rightarrow R^1}$; через $\{{\cal


D}(\Omega)\}^3$ --- пространство пробных функций


$\boldsymbol\varphi:\Omega\rightarrow R^3$ (${\boldsymbol{\varphi}=\{


\varphi_1,\varphi_2,\varphi_3\}}$,


$\varphi_i\in{\cal D}(\Omega)$, $i=1,2,3$);


${\cal D}(\overline{\Omega})$ --- множество функций, являющихся


сужениями на $\Omega$ функций из ${\cal D}(R^3)$;


$D'(\Omega)$ ---


пространство распределений (обобщенных функций) (\cite{3}, с.\,19).





Пусть $\{L_p(\Omega)\}^3$, $1\leq p<\infty$, --- банаховы


пространства суммируемых


со степенью $p$ вектор-функций $\bold u:\Omega\rightarrow R^3$ с


соответствующими нормами


$$


\|\bold{u}\|_{\{L_p(\Omega)\}^3}=\bigg\{\int_{\Omega}|\bold{u}


(\bold{x})|^pd\bold{x}\bigg\}^{1/p};


$$


здесь


$|\bold{u}(\bold{x})|=(u_1^2(\bold{x})+u_2^2(\bold{x})+u_3^2(\bold{x}))^{1/2}$.


При $p=2$ пространство $\{L_2(\Omega)\}^3$ является гильбертовым


со скалярным произведением


$(\bold{u},\bold{v})_{\{L_2{\Omega}\}^3}=\int_{\Omega}(\bold{u}


(\bold{x})\cdot\bold{v}(\bold{x}))d\bold{x}$.





Будем говорить, что $\Div\bold{u}=g\in L_p(\Omega)$ для 


функции $\bold{u}\in\{L_1(\Omega)\}^3$, если


$$


\int_{\Omega}g\varphi d\bold{x}=-\int_{\Omega}


(\Grad\varphi\cdot\bold{u})d\bold{x}\ \ \text{при всех}\ \ \varphi\in


{\cal D}(\Omega),\ \ \text{и}


$$


$\Rot\bold{u}=\bold{h}\in\{L_p(\Omega)\}^3$, если


$$


\int_{\Omega}(\bold{h}\cdot\boldsymbol{\psi})d\bold{x}=\int_{\Omega}


(\Rot\boldsymbol{\psi}\cdot\bold{u})d\bold{x}\ \ \text{при всех}


\ \ \boldsymbol{\psi}\in\{{\cal D}(\Omega)\}^3.


$$





Определим следующие пространства вектор-функций:


\begin{align*}


H_p(\Div;\Omega)&=\{\bold{u} \in \{ L_p(\Omega)\}^3\mid


\Div\bold{u}\in L_p(\Omega)\},\\


%


H_p(\Rot;\Omega)&=\{ \bold{u} \in


\{L_p(\Omega)\}^3 \mid \Rot\bold{u}\in \{L_p(\Omega)\}^3\}.


\end{align*}





\begin{lems}


$H_p(\Div;\Omega)$, $H_p(\Rot;\Omega)$ --- банаховы пространства с


соответствующими нормами


\begin{align*}


\|\bold{u}\|_{H_p(\Div;\Omega)}&=\{ \|\bold{u}


\|^p_{\{ L_p(\Omega)\}^3}+\|\Div\bold{u}\|^p_{L_p(\Omega)}\}^{1/p},\\


%%


\| \bold{u}


\|_{H_p(\Rot;\Omega)}&=\{ \| \bold{u} \|^p_{\{ L_p(\Omega)\}^3}+


\|\Rot \bold{u} \|^p_{\{ L_p(\Omega)\}^3}\}^{1/p}.


\end{align*}


\end{lems}





Через $H_p^0(\Div;\Omega)$ и $H_p^0(\Rot;\Omega)$


обозначим замыкание множества $\{{\cal D}(\Omega)\}^3$ в


$H_p(\Div;\Omega)$ и в $H_p(\Rot;\Omega)$ соответственно.





\begin{notenonum}


Если ограниченная область $\Omega$ имеет гладкую границу, то при $p=2$ 


можно показать, что


\begin{align*}


H_0(\Div;\Omega)&=H_2^0(\Div;\Omega)=\{\bold{u}\in


H(\Div;\Omega),\ \bold{u}\times\bold{n}=0,\ \bold{x}\in\Gamma\},\\


%


H_0(\Rot;\Omega)&=H_2^0(\Rot;\Omega)=\{\bold{u}\in


H(\Rot;\Omega),\ \bold{u}\cdot\bold{n}=0,\ \bold{x}\in\Gamma\},


\end{align*}


где граничные операторы понимаются в смысле теории следов (\cite{4}, 


с.\,19; \cite{5}, с.\,315).


\end{notenonum}


\addtocounter{section}{1}


\setcounter{equation}{0}


\setcounter{thms}{0}


\setcounter{lems}{0}





{\bf 2. Основные результаты.} Неравенства, приведенные в теоремах 2.1--2.3,


будут доказаны в п.\,4 на основании рассмотренных в п.\,3 представлений


вектор-функций в звездных областях.





Пусть $\Omega\subset


R^3$ --- открытое ограниченное множество с диаметром $d(\Omega)>0$, звездное


относительно некоторой точки.





\begin{thms}\label{t1}


При ${p>3/2}$, ${q=p/(p-1)}$  существует


положительная постоянная $C_1$, зависящая только от области $\Omega$ и $p$,


такая, что при любых $\bold{u}\in H_p(\Div;\Omega)$, ${\bold{v}\in


H_q^0(\Rot;\Omega)}$ будет справедливо неравенство


\begin{equation}


\bigg|\int_{\Omega}(\bold{u}\cdot


\bold{v})dx\bigg| \leq C_1\big(\|\bold{u}\|_{\{ L_p(\Omega)\}^3}\| \Rot


\bold{v} \|_{\{ L_q(\Omega)\}^3}+\| \bold{v} \|_{\{ L_q(\Omega)\}^3}\|


\Div\bold{u} \|_{L_p(\Omega)}\big).


\label{o1}


\end{equation}


\end{thms}








\begin{thms}\label{t2}


При ${p>3/2}$, ${q=p/(p-1)}$


существует положительная постоянная $C_2$, зависящая только от


области $\Omega$ и $p$, такая, что при любых $\bold{u}\in


H_q^0(\Div;\Omega)$, ${\bold{v}\in H_p(\Rot;\Omega)}$ будет


справедливо неравенство


\begin{multline}


\bigg|\int_{\Omega}(\bold{u}\cdot\bold{v})dx\bigg| \leq


C_2(\| \bold{u} \|_{\{ L_q(\Omega)\}^3}\| \Rot\bold{v} \|


_{\{ L_p(\Omega)\}^3}+ \\


%


+\| \bold{v} \|_{\{ L_p(\Omega)\}^3}\|


\Div\bold{u} \|_{L_q(\Omega)}+\|\Rot\bold{v}


\|_{\{ L_p(\Omega)\}^3}\|\Div\bold{u} \|_{L_q(\Omega)}).


\label{o2}


\end{multline}


\end{thms}








\begin{thms}\label{t3}


При ${p>3}$, ${q=p/(p-1)}$  существует


положительная постоянная $C_3$, зависящая только от области


$\Omega$ и $p$, такая, что при любых $\bold{u}\in H_q^0(\Div;\Omega)$,


$\bold{v}\in H_p(\Rot;\Omega)$ справедливо неравенство


\begin{equation}


\bigg|\int_{\Omega}(\bold{u}\cdot\bold{v})d\bold x\bigg|\leq


C_3(\| \bold{u} \|_{\{ L_q(\Omega)\}^3}\| \Rot\bold{v} \|


_{\{ L_p(\Omega)\}^3}+\| \bold{v} \|_{\{ L_p(\Omega)\}^3}\| 


\Div\bold{u} \|_{L_q(\Omega)}).


\label{o3} \end{equation}


\end{thms}


\addtocounter{section}{1}


\setcounter{equation}{0}


\setcounter{thms}{0}


\setcounter{lems}{0}





{\bf 3. Предварительные утверждения.}


Справедлива





\begin{thms}\label{tpl}


Пусть $\Omega\subset R^3$ ---


открытая, ограниченная, локально-звездная область. Тогда множество


вектор-функций, принадлежащих $\{{\cal D}(\overline{\Omega})\}^3$ плотно в


пространствах $H_p(\Div;\Omega)$ и $H_p(\Rot;\Omega)$.


\end{thms}





Доказательство этой теоремы для $p=2$


(\cite{4}, с.\,14; \cite{5}, с.\,314)


практически без изменений переносится на случай $1\leq p<\infty$.





Пусть область $\Omega\subset R^3$ звездна относительно


некоторой точки $\bold{y}\in\Omega$. Непосредственно проверяется





\begin{thms}


Для всех $\bold{u}\in\{C^1(\Omega)\}^3$ справедливы тождества


\begin{align}


\label{1}


\bold{u}(\bold{x})&=


\Grad_{\bold{x}}\bigg\{\int_0^1\bold{u}(\bold{z}_{\tau}


(\bold{x},\bold{y}))\cdot (\bold{x}-\bold{y})d\tau\bigg\}+\int_0^1\tau


[\Rot_{\bold{z}}\bold{u}(\bold{z}_{\tau}


(\bold{x},\bold{y}))\times (\bold{x}-\bold{y})]d\tau, \\


%


\label{2}


\bold{u}(\bold{x})&=\Rot_{\bold{x}}\bigg\{\int_0^1\tau[\bold{u}


(\bold{z}_{\tau}(\bold{x},\bold{y}))\times (\bold{x}-\bold{y})]d\tau\bigg\}


+\int_0^1\tau^2(\bold{x}-\bold{y})\Div_{\bold{z}}\bold{u}


(\bold{z}_{\tau}(\bold{x},\bold{y}))d\tau.


\end{align}


Здесь


$\bold{z}_{\tau}(\bold{x},\bold{y})=\bold{y}+\tau


(\bold{x}-\bold{y})\in\Omega$,


$\Rot_{\bold{z}}\bold{u}(\bold{z}_{\tau}(\bold{x},\bold{y}))=


\Rot_{\bold{z}}\bold{u}(\bold{z})$


в точке


${\bold{z}=\bold{z}_{\tau}(\bold{x},\bold{y})}$,


аналогично $\Div_{\bold{z}}\bold{u}(\bold{z}_


{\tau}(\bold{x},\bold{y}))=\Div_{\bold{z}}\bold{u}(\bold{z})$


в точке $\bold{z}=\bold{z}_{\tau}(\bold{x},\bold{y})$.


\end{thms}





Пусть $r=|\bold{x}-\bold{y}|$,


$\bold{s}=(\bold{x}-\bold{y})/r\in S=\{\boldsymbol{\theta}\in R^3$,


$|\boldsymbol{\theta}|=1\}$,


$\xi=\tau r$. Тогда тожества \eqref{1}, \eqref{2}


запишутся в виде


\begin{align}


\label{3}


\bold{u}(\bold{x})&=\Grad_{\bold{x}}\bigg\{\int_0^r(\bold{u}(


\bold{y}+\xi\bold{s})\cdot \bold{s})d\xi\bigg\}+\frac{1}{r}\int_0^r\xi


[\Rot\bold{u}(\bold{y}+\xi\bold{s})\times \bold{s}]d\xi, \\


%


\label{4}


\bold{u}(\bold{x})&=\Rot_{\bold{x}}\bigg\{\int_0^r\xi[\bold{u}(


\bold{y}+\xi\bold{s})\times


\bold{s}]d\xi\bigg\}+\frac{\bold{s}}{r^2}\int_0^r\xi^2


\Div\bold{u}(\bold{y}+\xi\bold{s})d\xi.


\end{align}





Получим некоторые оценки для интегралов, присутствующих в


\eqref{1}--\eqref{4}.





Пусть $\Omega \subset R^3 $ --- открытое


ограниченное множество с диаметром ${d(\Omega)>0}$, звездное относительно


некоторой точки $\bold{y}\in\Omega$. 


Для каждого $m>-1$ и функций $\varphi\in C(\overline{\Omega})$ и


$\boldsymbol{\varphi}\in \{C(\overline{\Omega})\}^3$ определим функции


$Q_m(\varphi):\Omega\rightarrow R^1$,


$\bold{Q}_m(\boldsymbol{\varphi}):\Omega\rightarrow R^3$ соотношениями


\begin{align}


Q_m(\varphi)(\bold{x})&=r^{-m}\int_0^r\xi^m\varphi(\bold{y}+


\xi\bold{s})d\xi,\label{5.1}\\


\bold{Q}_m(\boldsymbol{\varphi})(\bold{x})&=r^{-m}\int_0^r\xi^m\boldsymbol{\varphi}


(\bold{y}+\xi\bold{s})d\xi.


\label{5}


\end{align}





\begin{lems}\label{l1}


При всех $m>-1$, $q>\max\{1,3/(m+1)\}$ существует


такая положительная постоянная $C_m(q,\Omega)>0$, зависящая только от $m$,


$q$ и области $\Omega$, что для любых $\varphi\in C(\overline{\Omega})$


справедливо неравенство


\begin{equation}


\|Q_m(\varphi)\|_{L_q(\Omega)}\leq


C_m(q,\Omega)\|\varphi\|_{L_q(\Omega)}.


\label{oq}


\end{equation}


\end{lems}





{\bf Доказательство.} Имеем


\begin{multline*}


|Q_m(\varphi)(\bold{x})|\leq


r^{-m}\int_0^r\xi^m|\varphi


(\bold{y}+\xi\bold{s})|d\xi=r^{-m}\int_0^r\xi^{m-\frac{2}{q}}


\xi^{\frac{2}{q}}|\varphi(\bold{y}+\xi\bold{s})|d\xi\leq\\


%


\leq r^{-m}\bigg\{\int_0^r\xi^{(m-\frac{2}{q})\frac{q}{q-1}}d\xi\bigg\}^


{\frac{q-1}{q}}\bigg\{\int_0^r\xi^2|\varphi(\bold{y}+\xi\bold{s})|^qd\xi


\bigg\}^{1/q}= \\


%


=\bigg(\frac{q-1}{q(m+1)-3}\bigg)^{\frac{q-1}{q}}r^{\frac{q-3}{q}}


\bigg\{\int_0^r\xi^2|\varphi(\bold{y}+\xi\bold{s})|^qd\xi\bigg\}^{1/q}.


\end{multline*}





Пусть, далее, $R_{\bold{s}}(\bold{y})=\sup\{r\in


R^1:\bold{y}+r\bold{s}\in\Omega\}$


(${0<R_{\bold{s}}(\bold{y})<d(\Omega)}$), $d\bold{s}$ --- элемент площади


единичной сферы $S$. Тогда получим


\begin{gather*}


\|Q_m(\varphi)\|_{L_q(\Omega)}^q=\int_Sd\bold{s}


\int_0^{R_{\bold{s}}(\bold{y})}r^2|Q_m(\varphi)(\bold{y}+r\bold{s})|^q


dr\leq \\


%


\leq A\int_S


d\bold{s}\int_0^{R_{\bold{s}}(\bold{y})}r^{q-1}\int_0^r\xi^2


|\varphi(\bold{y}+\xi\bold{s})|^qd\xi dr\leq \\


%


\leq A \int_Sd\bold{s}


\int_0^{R_{\bold{s}}(\bold{y})}r^{q-1}\int_0^


{R_{\bold{s}}(\bold{y})}\xi^2|\varphi(\bold{y}+\xi\bold{s})|^qd\xi


dr\leq \\


%


\leq A \int_Sd\bold{s}


\frac{R_{\bold{s}}(\bold{y})}{q}\int_0^


{R_{\bold{s}}(\bold{y})}\xi^2|\varphi(\bold{y}+\xi\bold{s})|^qd\xi


dr\leq \\


%


\leq A \frac{(d(\Omega))^q}


{q}\int_Sd\bold{s}\int_0^{R_{\bold{s}}(\bold{y})}\xi^2


|\varphi(\bold{y}+\xi\bold{s})|^qd\xi dr=


A\bigg(\frac{d(\Omega)}{q^{1/q}}\|\varphi\|_{L_q(\Omega)}\bigg)^q,\\


%


A=\bigg(\frac{q-1}{q(m+1)-3}\bigg)^{q-1}.


\end{gather*}


Остается положить


\begin{equation}


\label{cq}


C_m(q,\Omega)=A^{{1}/{q}}


\frac{d(\Omega)}{q^{1/q}}.\quad\square


\end{equation}


В частности, из леммы \ref{l1} следуют справедливые при всех


$\varphi\in C(\overline{\Omega})$ оценки


\begin{alignat*}{2}


\|Q_0(\varphi)\|_{L_q(\Omega)}&\leq\bigg(\frac{q-1}{q-3}\bigg)


^{\frac{q-1}{q}}\frac{d(\Omega)}{q^{1/q}}


\|\varphi\|_{L_q(\Omega)}, &\quad q&>3,\\


%


\|Q_1(\varphi)\|_{L_q(\Omega)}&\leq\bigg(\frac{q-1}{2q-3}\bigg)


^{\frac{q-1}{q}}\frac{d(\Omega)}{q^{1/q}}


\|\varphi\|_{L_q(\Omega)}, &\quad q&>\frac{3}{2},\\


%


\|Q_2(\varphi)\|_{L_q(\Omega)}&\leq\frac{3^{1/q}


d(\Omega)}{3q^{1/q}}


\|\varphi\|_{L_q(\Omega)}, &\quad q&>1.


\end{alignat*}


Для функции $\boldsymbol{\varphi}\in\{ C(\overline{\Omega})\}^3$, применяя


лемму \ref{l1} к функции ${|\boldsymbol{\varphi}|\in C(\overline{\Omega})}$, имеем


$$


\|\bold{Q}_m(\boldsymbol{\varphi})\|_{\{L_q(\Omega)\}^3}\leq


\|Q_m(|\boldsymbol{\varphi}|)\|_{L_q(\Omega)}\leq C_m(q,\Omega)


\||\boldsymbol{\varphi}|\|_{L_q(\Omega)}=C_m(q,\Omega)


\|\boldsymbol{\varphi}\|_{\{L_q(\Omega)\}^3}.


$$





Полученные оценки показывают, что оператор $Q_m$, рассматриваемый как


отображение из $L_q(\Omega)$ в $L_q(\Omega)$, непрерывен на множестве


$C(\overline{\Omega})$, плотном в $L_q(\Omega)$.


Следовательно, его можно продолжить до


некоторого линейного ограниченного оператора, обозначаемого в дальнейшем


также через $Q_m$, определенного на всем пространстве $L_q(\Omega)$.


Оценка \eqref{oq} остается справедливой и для этого оператора.





Аналогично, оператор $\bold{Q}_m$ продолжается по непрерывности


на $\{L_q(\Omega)\}^3$.





\begin{thms} \label{n1}


Пусть $\Omega\subset R^3$ --- открытое ограниченное множество, 


звездное относительно некоторой точки $y\in\Omega$. Для всех


$\bold{u}\in H_p(\Div;\Omega)$ при ${p>3/2}$ справедливо тождество


\begin{equation}


\bold{u}(\bold{x})=


\Rot_{\bold{x}}[\bold{Q}_1(\bold{u})\times\bold{s}]+\bold{s}Q_2(\Div\bold{u}).


\label{nn1}


\end{equation}


\end{thms}





\begin{pf}


Пусть $\bold{u}\in H_p(\Div;\Omega)$,


$p>3/2$. Согласно теореме \ref{tpl} найдется последовательность


$\{\bold{u}_n\}\subset\{{\cal D}(\overline{\Omega})\}^3$ такая, что


$\bold{u}_n\to\bold{u}$ при $n\to\infty$ в $H_p(\Div;\Omega)$.


Для всех $\bold{u}_n$ справедливо представление


\eqref{3}, которое с учетом \eqref{5.1}, \eqref{5}


можно записать в виде


\begin{equation}


\bold{u}_n(\bold{x})=\Rot_{\bold{x}}[\bold{Q}_1(\bold{u}_n)\times\bold{s}]+


\bold{s}Q_2(\Div\bold{u}_n).


\label{b1}


\end{equation}


Ввиду \eqref{oq}


\ $\bold{u}_n\to\bold{u}$,


$\bold{s}Q_2(\Div\bold{u}_n)\to\bold{s}Q_2(\Div\bold{u})$


в $\{L_p(\Omega)\}^3$ при $n\to\infty$.





Последовательность


$\{\Rot_{\bold{x}}[\bold{Q}_1(\bold{u}_n)\times\bold{s}]\}$, таким образом,


фундаментальна в $\{L_p(\Omega)\}^3$,


следовательно, сходится к некоторому элементу $\bold{v}\in\{L_p(\Omega)\}^3$.


С другой стороны, $\Rot_{\bold{x}}[\bold{Q}_1(\bold{u}_n)\times\bold{s}]


\to\Rot_{\bold{x}}[\bold{Q}_1(\bold{u})\times\bold{s}]$


в пространстве распределений. Действительно,


\begin{multline*}


|\la \Rot_{\bold{x}}[\bold{Q}_1(\bold{u}_n)\times\bold{s}]


-\Rot_{\bold{x}}[\bold{Q}_1(\bold{u})\times\bold{s}],\boldsymbol{\psi}\ra|=


|\la[\bold{Q}_1(\bold{u}_n-\bold{u})\times\bold{s}],\Rot\boldsymbol{\psi}\ra|\leq \\


%


\leq C_1(p,\Omega)\|\bold{u}_n-\bold{u}\|_{\{L_p(\Omega)\}^3}\|\Rot


\boldsymbol{\psi}\|_{\{L_p(\Omega)\}^3}\to 0


\end{multline*}


для любой $\boldsymbol{\psi}\in\{{\cal D}(\Omega)\}^3$.


Так как функции


$\Rot_{\bold{x}}[\bold{Q}_1(\bold{u})\times\bold{s}]$ и $\bold{v}$


определяют одинаковые распределения, то $\Rot_{\bold{x}}[\bold{Q}_1


(\bold{u})\times\bold{s}]=\bold{v}$, и, переходя к пределу в равенстве \eqref{b1},


получим \eqref{nn1}.


\end{pf}





\begin{thms}\label{n2}


Пусть $\Omega\subset R^3$ --- открытое ограниченное множество,


звездное относительно некоторой точки.


Для любой $\bold{u}\in H_p(\Rot;\Omega)$ при


${p>3/2}$ найдется функция $Q\in L_p(\Omega)$ такая, что справедливо


тождество


\begin{equation}


\bold{u}(\bold{x})=\Grad_{\bold{x}}Q+[\bold{Q}_1(\Rot


\bold{u})\times\bold{s}],


\label{nn2}


\end{equation}


причем, если $p>3$, можно взять $Q=(\bold{Q}_0(\bold{u})\cdot\bold{s})$.


Дифференциальные операторы понимаются в смысле теории распределений.


\end{thms}





\begin{pf}


Пусть $\bold{u}\in H_p(\Rot;\Omega)$,


$p>3/2$. Согласно теореме \ref{tpl} найдется последовательность


$\{\bold{u}_n\}\subset\{{\cal D}(\overline{\Omega})\}^3$ такая, что


$\bold{u}_n\to\bold{u}$ при $n\to\infty$ в $H_p(\Rot;\Omega)$.


Для всех $\bold{u}_n$ справедливо представление


\eqref{4}, которое с учетом \eqref{5.1}, \eqref{5} можно записать в виде


\begin{equation}


\bold{u}_n(\bold{x})=\Grad_{\bold{x}}(\bold{Q}_0(\bold{u}_n)\cdot\bold{s})+


[\bold{Q}_1(\Rot \bold{u}_n)\times\bold{s}].


\label{b2}


\end{equation}


При $n\to\infty$ \ $\bold{u}_n\to\bold{u}$,


$[\bold{Q}_1(\Rot\bold{u}_n)\times\bold{s}]\to[\bold{Q}_1(


\Rot\bold{u})\times\bold{s}]$ в $\{L_p(\Omega)\}^3$ ввиду \eqref{oq}.





Последовательность


$\{\Grad_{\bold{x}}(\bold{Q}_0(\bold{u}_n)\cdot\bold{s})\}$, таким образом,


фундаментальна в $\{L_p(\Omega)\}^3$,


следовательно, сходится к некоторому $\bold{v}\in\{L_p(\Omega)\}^3$.


Известны следующие утверждения.





\begin{lems}[{\series{m}\shape{n}\selectfont \cite{4}, с.\,20) (характеризация градиента распределения}]


Пусть ${\Omega\subset R^3}$ --- открытое множество,


$\bold{f}\in\{D'(\Omega)\}^3$,


$V=\{\bold{v}\in \{D(\Omega)\}^3;$


$\Div\bold{v}=0\}$. Для того чтобы $\bold{f}=\Grad p$ для некоторого $p\in


D'(\Omega)$, необходимо и достаточно, чтобы $\la\bold{f},\bold{v}\ra=0$ для всех


$\bold{v}\in V$.


\end{lems}








\begin{lems}[{\series{m}\shape{n}\selectfont \cite{9}, с.\,26) (неравенство Пуанкаре}]


\label{np}


Пусть $\Omega\subset R^3$ --- область класса $C^{0,1}$, $p>1$.


Существует положительная постоянная $A(\Omega,p)$, при которой для


каждой $u\in D'(\Omega)$ такой, что $\Grad u\in\{L_p(\Omega)\}^3$,


найдется такое число ${C^*>0}$, что


$$


\|u-C^*\|_{L_p(\Omega)}\leq A(\Omega,p)\|


\Grad u\|_{\{L_p(\Omega)\}^3}.


$$


\end{lems}





Следовательно, $\bold{v}=\Grad Q$, где $Q\in


L_p(\Omega)$ определена с точностью до аддитивной константы.





Если $p>3$, то


$\Grad_{\bold{x}}(\bold{Q}_0(\bold{u}_n)\cdot\bold{s})$ стремится


к $\Grad_{\bold{x}}(\bold{Q}_0(\bold{u})\cdot\bold{s})$ в


пространстве распределений, поэтому можно положить


$Q=(\bold{Q}_0(\bold{u})\cdot\bold{s})$.


Переходя к пределу в равенстве \eqref{b2}, получим \eqref{nn2}.


\end{pf}





Следствием теорем \ref{n1}, \ref{n2} является





\begin{lems}


Пусть $\Omega\subset R^3 $ --- открытое ограниченное множество,


звездное относительно некоторой точки. Для всех


${\bold{u}\in H_p(\Div;\Omega)}$


\ $[\bold{Q}_1(\bold{u})\times\bold{s}]\in H_p(\Rot;\Omega)$. Для всех


$\bold{u}\in H_p(\Rot;\Omega)$ функция $Q$,


определяемая теоремой $\ref{n2}$, лежит в пространстве Соболева


$W^{1,p}(\Omega)$.


\end{lems}


\addtocounter{section}{1}


\setcounter{equation}{0}


\setcounter{thms}{0}





{\bf 4. Доказательство основных неравенств.}


Для доказательства теоремы \ref{t1} положим\linebreak $\bold{u}\in\{{\cal


D}(\overline{\Omega})\}^3$, $\bold{v}\in\{{\cal D}(\Omega)\}^3$. Используя


представление \ref{nn1} для $\bold{u}$, получим


\begin{gather*}


\int_{\Omega}(\bold{u}(\bold{x})\cdot\bold{v}(\bold{x}))d\bold{x}=


\int_{\Omega}(\bold{v}(\bold{x})\cdot\Rot_{\bold{x}}[\bold{Q}_1


(\bold{u})\times\bold{s}])d\bold{x}+\int_{\Omega}(\bold{v}\cdot\bold{s})


Q_2(\Div\bold{u})d\bold{x}= \\


%


=\int_{\Omega}(\Rot\bold{v}(\bold{x})\cdot[\bold{Q}_1


(\bold{u})\times\bold{s}])d\bold{x}+\int_{\Omega}(\bold{v}\cdot\bold{s})


Q_2(\Div\bold{u})d\bold{x}+\int_{\Gamma}([\bold{v}(\bold{x})\times


[\bold{Q}_1(\bold{u})\times\bold{s}]]\cdot\bold{n}(\bold{x}))d\Gamma.


\end{gather*}


Последний интеграл, получившийся в результате применения тождества


$$


(\Rot\bold{a}\cdot\bold{b})=(\Rot\bold{b}\cdot\bold{a})+


\Div[\bold{a}\times\bold{b}]


$$


и теоремы Гаусса--Остроградского, равен нулю за счет


выбора $\bold{v}$. Используя неравенство Г\"ель\-де\-ра с показателями $p$ и $q$ и


оценки для $Q_m$, получим


\begin{multline*}


\bigg|\int_{\Omega}(\bold{u}(\bold{x})\cdot\bold{v}(\bold{x}))d\bold{x}\bigg|\leq


\|\Rot\bold{v}\|_{\{L_q(\Omega)\}^3}\|\bold{Q}_1(\bold{u})\|_


{\{L_p(\Omega)\}^3}+ 


\|\bold{v}\|_{\{L_q(\Omega)\}^3}\|Q_2(


\Div\bold{u})\|_ {L_p(\Omega)}\leq \\


%


\leq C_1(\|


\Rot\bold{v}\|_{\{L_q(\Omega)\}^3}\|\bold{u}\|_


{\{L_p(\Omega)\}^3}+\|\bold{v}\|_{\{L_q(\Omega)\}^3}\|\Div\bold{u}\|_


{L_p(\Omega)}),


\end{multline*}


где $C_1=C_1(\Omega,p)$.





Согласно теореме о плотности и определению


пространства $H_q^0(\Rot;\Omega)$ найденная оценка справедлива


для $\bold{u}\in H_p(\Div;\Omega)$, $\bold{v}\in


H_q^0(\Rot;\Omega)$ и теорема \ref{t1} доказана.





Для доказательства теорем \ref{t2}, \ref{t3} возьмем $\bold{v}\in\{{\cal


D}(\overline{\Omega})\}^3$, $\bold{u}\in\{{\cal D}(\Omega)\}^3$. Из


представления \eqref{nn2} для $\bold{v}$ и оценки для $Q_1$


получаем


$$\|\Grad_{\bold{x}}(\bold{Q}_0(\bold{v})\cdot\bold{s})\|_


{\{L_p(\Omega)\}^3}\leq\|\bold{v}\|_{\{L_p(\Omega)\}^3}+


\|\bold{Q}_1(\Rot\bold{v})\|_{\{L_p(\Omega)\}^3}\leq


\|\bold{v}\|_{\{L_p(\Omega)\}^3}+C_1(\Omega,p)


\|\Rot\bold{v}\|_{\{L_p(\Omega)\}^3}.


$$


Согласно лемме \ref{np} найдется такое число $C^*>0$, зависящее от $\bold{v}$,


что


\begin{multline}


\|(\bold{Q}_0(\bold{v})\cdot\bold{s})-C^*\|_{L_p(\Omega)}\leq


A(\Omega,p)\|\Grad_{\bold{x}}(\bold{Q}_0(\bold{v})\cdot\bold{s})\|_


{\{L_p(\Omega)\}^3}\leq \\


%


\leq A(\Omega,p)\|\bold{v}\|_{\{L_p(\Omega)\}^3}+A(\Omega,p)C_1(\Omega,p)


\|\Rot\bold{v}\|_{\{L_p(\Omega)\}^3}.


\label{30}


\end{multline}


Далее,


\begin{gather*}


\int_{\Omega}(\bold{u}(\bold{x})\cdot\bold{v}(\bold{x}))d\bold{x}=


\int_{\Omega}(\bold{u}(\bold{x})\cdot\Grad_{\bold{x}}\{(\bold{Q}_0


(\bold{u})\cdot\bold{s})-C^*\})d\bold{x}+ \\


%


+\int_{\Omega}


(\bold{u}\cdot[\bold{Q}_1(\Rot\bold{u})\times\bold{s}])d\bold{x}=


-\int_{\Omega}((\bold{Q}_0(\bold{u})\cdot\bold{s})-C^*)


\Div\bold{u}d\bold{x}+ \\


%


+\int_{\Omega}


(\bold{u}\cdot[\bold{Q}_1(\Rot\bold{u})\times\bold{s}])d\bold{x}+


\int_{\Gamma}((\bold{Q}_0(\bold{u})\cdot\bold{s})-С^*)


\bold{u}_nd\Gamma.


\end{gather*}


Последний интеграл, полученный в результате применения тождества


$$


b\Div\bold{a}=\Div b\bold{a}-(\Grad b\cdot\bold{a})


$$


и теорeмы Гаусса--Остроградского, равен нулю за счет


выбора $\bold{u}$. Применяя неравенство Г\"ель\-де\-ра с показателями $p$ и $q$,


оценки для $Q_m$ и неравенство \ref{30}, получим


\begin{gather*}


\bigg|\int_{\Omega}(\bold{u}(\bold{x})\cdot\bold{v}(\bold{x}))d\bold{x}\bigg|\leq


\|\Div\bold{u}\|_{L_q(\Omega)}(A(\Omega,p)\|\bold{v}\|_


{\{L_p(\Omega)\}^3}+ \\


%


+A(\Omega,p)C_1(\Omega,p)


\|\Rot\bold{v}\|_{\{L_p(\Omega)\}^3})+


\|\bold{u}\|_{\{L_q(\Omega)\}^3}\|\bold{Q}_1(\Rot\bold{u})\|_


{L_p(\Omega)}\leq \\


%


\leq C_2(\|\Div\bold{u}\|_{L_q(\Omega)}\|\bold{v}\|_


{\{L_p(\Omega)\}^3}+\|\Div\bold{u}\|_{L_q(\Omega)}\|\Rot\bold{v}\|_


{\{L_p(\Omega)\}^3}+


\|\bold{u}\|_{\{L_q(\Omega)\}^3}\|\Rot\bold{v}\|_


{\{L_p(\Omega)\}^3}),


\end{gather*}


где $C_2=\max(A(\Omega,p),C_1(\Omega,p),A(\Omega,p)C_1(\Omega,p))$.





C помощью предельного перехода, теоремы


о плотности и определения пространства $H_q^0(\Div;\Omega)$ установим


справедливость оценки \eqref{o2} для $\bold{u}\in H_q^0(\Div;\Omega)$,


${\bold{v}\in H_p(\Rot;\Omega)}$. Теорема \ref{t2}, таким образом, доказана.





Если $p>3$, из представления \ref{b2}


для $\bold{v}$ и оценок для $\bold{Q_1}$, $\bold{Q}_0$ имеем


\begin{gather*}


\int_{\Omega}(\bold{u}(\bold{x})\cdot\bold{v}(\bold{x}))d\bold{x}=


\int_{\Omega}(\bold{u}(\bold{x})\cdot\Grad_{\bold{x}}(\bold{Q}_0


(\bold{u})\cdot\bold{s}))d\bold{x}+\int_{\Omega}


(\bold{u}\cdot[\bold{Q}_1(\Rot\bold{u})\times\bold{s}])d\bold{x}=\\


%


=-\int_{\Omega}(\bold{Q}_0(\bold{u})\cdot\bold{s})


\Div\bold{u}d\bold{x}+\int_{\Omega}


(\bold{u}\cdot[\bold{Q}_1(\Rot\bold{u})\times\bold{s}])d\bold{x},\\


%


\bigg|\int_{\Omega}(\bold{u}(\bold{x})\cdot\bold{v}(\bold{x}))d\bold{x}\bigg|\leq


C_0(\Omega,p)\|\Div\bold{u}\|_{L_q(\Omega)}\|\bold{v}\|_


{\{L_p(\Omega)\}^3}+\\


%


+\|\bold{u}\|_{\{L_q(\Omega)\}^3}\|\bold{Q}_1(\Rot\bold{u})\|_


{\{L_p(\Omega)\}^3}\leq \\


%


\leq C_3(\|\Div\bold{u}\|_{L_q(\Omega)}\|\bold{v}\|_


{\{L_p(\Omega)\}^3}+\|\bold{u}\|_{\{L_q(\Omega)\}^3}\|\Rot\bold{v}\|_


{\{L_p(\Omega)\}^3}),


\end{gather*}


где $C_3=C_0(\Omega,p)$.





Так как $\{{\cal D}(\overline{\Omega})\}^3$ плотно в пространстве


$H_p(\Rot;\Omega)$, а множество $\{{\cal D}(\Omega)\}^3$ ---


в пространстве $H_q^0(\Div;\Omega)$,


теорема \ref{t3} доказана.


\addtocounter{section}{1}


\setcounter{equation}{0}


\setcounter{thms}{0}


\setcounter{lems}{0}


\medskip





{\bf 5. Некоторые следствия основных неравенств для случая $p=2$.}


Обозначим


$H(\Rot;\Omega)=H_2(\Rot;\Omega)$,


$H(\Div;\Omega)=H_2(\Div;\Omega)$.





Поскольку неравенства \eqref{o1}, \eqref{o2} при $p=q=2$ оценивают скалярные


произведения функций в $\{L_2(\Omega)\}^3$, из них следуют известные условия


ортогональности соленоидальных и потенциальных векторных полей в


$\{L_2(\Omega)\}^3$ (см. \cite{1}--\cite{8}) для звездных


областей.





Далее, из неравенств \eqref{o1}, \eqref{o2} при $\bold{v}=\bold{u}$


вытекают следствия теорем \ref{t1}, \ref{t2}.


Пусть $\Omega\subset


R^3$ --- открытое ограниченное множество с диаметром $d(\Omega)>0$, звездное


относительно некоторой точки. 





\begin{thms}\label{t6}


Существует положительная постоянная


$C_4$, зависящая только от области $\Omega$, такая, что при любых


$\bold{u}\in H(\Div;\Omega)\bigcap H^0(\Rot;\Omega)$ будет


выполняться неравенство


$$


\|\bold{u}\|_{\{L_2(\Omega)\}^3}\leq C_4(\|


\Rot\bold{u}\|_{\{L_2(\Omega)\}^3}+\|\Div \bold{u}\|_{L_2(\Omega)}).


$$


\end{thms}





\begin{thms}\label{t7}


Существует


положительная постоянная $C_5$, зависящая только от области $\Omega$,


такая, что при любых $\bold{u}\in H^0(\Div;\Omega)\bigcap H(


\Rot;\Omega)$ будет выполняться неравенство 


$$


\|\bold{u}\|_{\{L_2(\Omega)\}^3}\leq C_5(\|


\Rot\bold{u}\|_{\{L_2(\Omega)\}^3}+\|\Div \bold{u}\|_{L_2(\Omega)}).


$$


\end{thms}





\begin{pf}


Справедливо неравенство \eqref{o2}, где


$\bold{v}=\bold{u}$, т.\,е.


$$


\bigg|\int_{\Omega}\bold{u}^2dx\bigg|\leq C_2(\|


\bold{u} \|_{\{ L_2(\Omega)\}^3}(\| \Rot\bold{u} \|_{\{


L_2(\Omega)\}^3}+\| \Div\bold{u}


\|_{L_2(\Omega)})+\|\Rot\bold{u} \|_{\{ L_2(\Omega)\}


^3}\|\Div \bold{u} \|_{L_2(\Omega)}).


$$





Обозначив $z=\|\bold{u} \|_{\{


L_2(\Omega)\}^3}$, $a=\|\Div \bold{u} \|_{L_2(\Omega)}$,


$b=\|\Rot\bold{u} \|_{\{ L_2(\Omega)\}^3}$ и разрешив неравенство


$$


z^2-C_2(a+b)z-C_2ab\leq 0


$$


относительно $z\geq 0$, получим


\begin{multline*}


z\leq\frac{C_2(a+b)+\sqrt{C_2^2(a+b)^2+4C_2ab}}{2}\leq


\frac{C_2(a+b)+\sqrt{(a+b)^2(C_2^2+2C_2)}}{2}\leq\\


\leq(a+b)\frac{C_2+\sqrt{C^2_2+2C_2}}{2}\leq


(a+b) (C_2+\tfrac{1}{2}),


\end{multline*}


откуда и следует требуемая оценка с константой $C_5=C_2+1/2$.


\end{pf}





Пусть теперь $\Omega\subset R^3$ --- открытое ограниченное


множество (не обязательно звездное).





Обозначим через $\ker(\Div;\Omega)$,


$\ker(\Rot;\Omega)$


ядра операторов $\Div$ и $\Rot$ соответственно,


$\ker^{\bot}(\Div;\Omega)$, $\ker^{\bot}(\Rot;\Omega)$ ---


ортогональные дополнения к ним в $\{L_2(\Omega)\}^3$.





\begin{lems}\label{do}


$\ker^{\bot}(\Div;\Omega)$ совпадает с


замыканием в $\{L_2(\Omega)\}^3$ множества $\{


\Grad\varphi,\varphi\in{\cal D}(\Omega)\}$.


\end{lems}





\begin{pf}


Обозначим замыкание множества


$\{\Grad\varphi,\varphi\in{\cal D}(\Omega)\}$ в $\{L_2(\Omega)\}^3$


че\-рез~$M$. Пусть $\varphi\in{\cal D}(\Omega)$, $\bold{v}\in\ker(\Div;\Omega)$,


$\int\limits_{\Omega}(\bold{v}\cdot 


\Grad\varphi)d\bold{x}=0$ по определению $\Div\bold{v}$, т.\,е.


$\Grad\varphi\in\ker^{\bot}(\Div;\Omega)$. Так как $\ker^{\bot}(


\Div;\Omega)$ --- замкнутое подпространство $\{L_2(\Omega)\}^3$, то


$M\subset\ker^{\bot}(\Div;\Omega)$.





Обратно, пусть $M$ --- собственное подпространство в


$\ker^{\bot}(\Div;\Omega)$. Тогда найдется ненулевой элемент


$\bold{z}\in\ker^{\bot}(\Div;\Omega)$, ортогональный $M$. В частности,


$\int\limits_{\Omega}(\bold{z}\cdot 


\Grad\varphi)d\bold{x}=0$ для всех $\varphi\in{\cal D}(\Omega)$, т.\,е.


$\bold{z}\in\ker(\Div;\Omega)$, и, следовательно,


$\bold{z}=\bold{0}$. Полученное противоречие доказывает лемму.


\end{pf}





Аналогично доказывается





\begin{lems}\label{ro}


$\ker^{\bot}(\Rot;\Omega)$ совпадает с


замыканием в $\{L_2(\Omega)\}^3$ множества $\{


\Rot\boldsymbol{\psi},\boldsymbol{\psi}\in\{{\cal D}(\Omega)\}^3\}$.


\end{lems}





\begin{thms}[{\series{m}\shape{n}\selectfont теорема Юнга}]


Каждое открытое множество


${\Omega\subset R^n}$ с диаметром $d(\Omega){>}0$


принадлежит некоторому открытому шару


${B_R(\bold{y})=\{\bold{x}\in


R^n:|\bold{x}-\bold{y}|\leq R\}}$, где


\begin{equation}


R=\sqrt{\frac{n}{2(n+1)}}d(\Omega). \label{un}


\end{equation}


\end{thms}





\begin{thms}


Пусть $\Omega\subset R^3$ --- открытое ограниченное множество с


диаметром $d(\Omega)>0$. Тогда при всех $\bold{u}\in


H_0(\Div;\Omega)\bigcap \ker^{\bot}(\Div;\Omega)$ выполняется


неравенство


$$


\int_{\Omega}\bold{u}^2(\bold{x})d\bold{x}\leq\frac{d^2(\Omega)}{16}


\int_{\Omega}|\Div\bold{u}(\bold{x})|^2d\bold{x}.


$$


\end{thms}





\begin{pf}


Пусть $\bold{u}\in


H_0(\Div;\Omega)\bigcap \ker^{\bot}(\Div;\Omega)$. Тогда существуют


последовательности $\{\varphi^n\}\in {\cal D}(\Omega)$, $\{\boldsymbol{\psi}^n\}\in


\{{\cal D}(\Omega)\}^3$ пробных функций такие, что


$\|\bold{u}-\Grad\varphi^n\|_{\{L_2(\Omega)\}^3}\to 0$,


$\|\bold{u}-\boldsymbol{\psi}^n\|_{H(\Div;\Omega)}\to 0$


при $n\to\infty$.





По теореме Юнга найдется $\bold{y}\in R^3$, при котором $\Omega\subset


B_R(\bold{y})$, где $R$ определяется соотношением \eqref{un}.





Определим для каждого $n=1,2,\dotsc$ пробные 


функции $\varphi_R^n\in {\cal D}(B_R(\bold{y}))$,


${\boldsymbol{\psi}_R^n\in\{{\cal D}(B_R(\bold{y}))\}^3}$


соотношениями


$$


\varphi_R^n=


\begin{cases}


\varphi^n, & \text{если }\ \bold{x}\in\Omega;\\


0, & \text{если }\ \bold{x}\in B_R(\bold{y})\setminus \Omega,


\end{cases}\qquad


\boldsymbol{\psi}_R^n=


\begin{cases}


\boldsymbol{\psi}^n, & \text{если }\ \bold{x}\in\Omega;\\


0, & \text{если }\ \bold{x}\in


B_R(\bold{y})\setminus \Omega.


\end{cases}


$$





Согласно \eqref{2} при всех $\bold{x}\in B_R(\bold{y})$ имеем


$$


\Div\boldsymbol{\psi}_R^n(\bold{x})=\Div_{\bold{x}}\bigg\{\int_0^1


\tau^2(\bold{x}-\bold{y})


\Div\boldsymbol{\psi}_R^n(\tau\bold{x}+(1-\tau)\bold{y})d\tau\bigg\}.


$$





Умножая это


равенство на $\varphi_R^n$ и интегрируя по шару $B_R(\bold{y})$, получим


$$


\int\limits_{B_R(\bold{y})}\varphi_R^n(\bold{x})\Div\boldsymbol{\psi}_R^n


(\bold{x})d\bold{x}=


\int\limits_{B_R(\bold{y})}\varphi_R^n(\bold{x})


\Div_{\bold{x}}\bigg\{\int_0^1\tau^2(\bold{x}-\bold{y})\Div\boldsymbol{\psi}_R^n


(\tau\bold{x}+(1-\tau)\bold{y})d\tau\bigg\}d\bold{x}


$$


или


\begin{gather*}


\bigg|\int\limits_{B_R(\bold{y})}(\boldsymbol{\psi}_R^n(\bold{x})\cdot\Grad


\varphi_R^n(\bold{x}))d\bold{x}\bigg|=\\


%


=\bigg|\int\limits_{B_R(\bold{y})}(\Grad


\varphi_R^n(\bold{x})\cdot\int_0^1\tau^2(\bold{x}-\bold{y})\Div\boldsymbol{\psi}_R^n


(\tau\bold{x}+(1-\tau)\bold{y})d\tau)d\bold{x}\bigg|\leq\\


%


\bigg\{\int\limits_{B_R(\bold{y})}(\Grad\varphi_R^n(\bold{x}))^2d\bold


{x}\bigg\}^{1/2}\bigg\{\int\limits_{B_R(\bold{y})}\bigg(\int_0^1\tau^2


(\bold{x}-\bold{y})\Div\boldsymbol{\psi}_R^n


(\tau\bold{x}+(1-\tau)\bold{y})d\tau\bigg)^2d\bold{x}\bigg\}^{1/2}.


\end{gather*}


В силу \eqref{oq}


$$


\bigg|\int\limits_{B_R(\bold{y})}(\boldsymbol{\psi}_R^n(\bold{x})\cdot


\Grad\varphi_R^n(\bold{x}))d\bold{x}\bigg|\leq 


\frac{R}{\sqrt{6}}\bigg\{\int\limits_{B_R (\bold{y})}(


\Grad\varphi_R^n(\bold{x}))^2d\bold


{x}\bigg\}^{1/2}\bigg\{\int\limits_{B_R(\bold{y})}(\Div


\boldsymbol{\psi}_R^n (\bold{x}))^2d\bold{x}\bigg\}^{1/2}


$$


или


$$


\bigg|\int_{\Omega}(\boldsymbol{\psi}_R^n(\bold{x})\cdot\Grad


\varphi_R^n(\bold{x}))d\bold{x}\bigg|\leq


\frac{R}{\sqrt{6}}\bigg\{\int_


{\Omega}(\Grad\varphi_R^n(\bold{x}))^2d\bold


{x}\bigg\}^{1/2}\bigg\{\int_{\Omega}(\Div\boldsymbol{\psi}_R^n


(\bold{x}))^2d\bold{x}\bigg\}^{1/2}.


$$


Переходя в последнем неравенстве к пределу при $n\to\infty$


и принимая во внимание \eqref{un}, получим


неравенство


$$


\int_{\Omega}\bold{u}^2(\bold{x})d\bold{x}\leq\frac{d(\Omega)}{4}


\bigg\{\int_{\Omega}\bold{u}^2(\bold{x})d\bold{x}\bigg\}^{1/2}


\bigg\{\int_{\Omega}|\Div\bold{u}(\bold{x})|^2d\bold{x}\bigg\}^{1/2},


$$


из которого следует требуемая оценка.


\end{pf}





\begin{thms}


Пусть $\Omega\subset R^3$ ---


открытое ограниченное множество с диаметром $d(\Omega)>0$. Тогда при всех


$\bold{u}\in H_0(\Rot;\Omega)\bigcap \ker^{\bot}(\Rot;\Omega)$


выполняется неравенство


$$


\int_{\Omega}\bold{u}^2(\bold{x})d\bold{x}\leq\frac{3d^2(\Omega)}{16}


\int_{\Omega}|\Rot\bold{u}(\bold{x})|^2d\bold{x}.


$$


\end{thms}





\begin{pf}


Пусть $\bold{u}\in H_0(\Rot;\Omega)\bigcap


\ker^{\bot}(\Rot;\Omega)$. Тогда существуют


последовательности $\{\boldsymbol{\varphi}^n\}\in \{{\cal D}(\Omega)\}^3$,


$\{\boldsymbol{\psi}^n\}\in \{{\cal D}(\Omega)\}^3$ пробных функций такие, что


$\|\bold{u}-\Rot\boldsymbol{\varphi}^n\|_{\{L_2(\Omega)\}^3}\to 0$,


$\|\bold{u}-\boldsymbol{\psi}^n\|_{H(\Rot;\Omega)}\to 0$


при $n\to\infty$.





По теореме Юнга найдется $\bold{y}\in R^3$, при


котором $\Omega\subset B_R(\bold{y})$, где


${R=\sqrt{3}d(\Omega)/2\sqrt{2}}$.





Определим для каждого $n=1,2,\dotsc$


пробные функции $\boldsymbol{\varphi}_R^n\in \{{\cal D}(B_R(\bold{y}))\}^3$,


$\boldsymbol{\psi}_R^n\in\{{\cal D}(B_R(\bold{y}))\}^3$ соотношениями


$$


\boldsymbol{\varphi}_R^n=


\begin{cases}


\boldsymbol{\varphi}^n, & \text{если }\ \bold{x}\in\Omega;\\


0, & \text{если }\ \bold{x}\in B_R(\bold{y})\setminus \Omega,


\end{cases}\qquad


\boldsymbol{\psi}_R^n=


\begin{cases}


\boldsymbol{\psi}^n, &\text{если }\ \bold{x}\in\Omega;\\


0, & \text{если \ }\bold{x}\in B_R(\bold{y})\setminus \Omega.


\end{cases}


$$





Согласно \eqref{1} при всех $\bold{x}\in B_R(\bold{y})$ имеем


$$


\Rot\boldsymbol{\psi}_R^n(\bold{x})=\Rot_{\bold{x}}\bigg\{\int_0^1\tau[


\Rot\boldsymbol{\psi}_R^n(\tau\bold{x}+(1-\tau)\bold{y})


\times(\bold{x}-\bold{y})]d\tau\bigg\}.


$$


Умножая это равенство скалярно на $\boldsymbol{\varphi}_R^n$ и интегрируя по шару


$B_R(\bold{y})$, получим


\begin{gather*}


\bigg|\int\limits_{B_R(\bold{y})}(\boldsymbol{\varphi}_R^n(\bold{x})\cdot\Rot


\boldsymbol{\psi}_R^n(\bold{x}))d\bold{x}\bigg|= \\


%


=\bigg|\int\limits_{B_R(\bold{y})}\bigg(\boldsymbol


{\varphi}_R^n(\bold{x})\cdot\Rot_{\bold{x}}\bigg\{\int_0^1\tau[


\Rot\boldsymbol{\psi}_R^n(\tau\bold{x}+(1-\tau)\bold{y})\times(\bold{x}-\bold{y})]


d\tau\bigg\}\bigg)d\bold{x}\bigg|= \\


%


=\bigg|\int\limits_{B_R(\bold{y})}\bigg(\Rot\boldsymbol


{\varphi}_R^n(\bold{x})\cdot\Rot\bigg\{\int_0^1


\tau[\Rot\boldsymbol{\psi}_R^n(\tau\bold{x}+(1-\tau)\bold{y})\times


(\bold{x}-\bold{y})]d\tau\bigg\}\bigg)d\bold{x}\bigg|\leq \\


%


\leq\bigg\{\int\limits_{B_R(\bold{y})}(\Rot\boldsymbol


{\varphi}_R^n(\bold{x}))^2d\bold{x}\bigg\}^{1/2}


\bigg\{\int\limits_{B_R (\bold{y})}\bigg(\int_0^1


\tau[\Rot\boldsymbol{\psi}_R^n(\tau\bold{x}+(1-\tau)\bold{y})\times


(\bold{x}-\bold{y})]d\tau\bigg)^2d\bold{x}\bigg\}^{1/2}.


\end{gather*}





В силу \eqref{cq}


$$


\bigg|\int\limits_{B_R(\bold{y})}(\boldsymbol{\varphi}_R^n(\bold{x})\cdot\Rot


\boldsymbol{\psi}_R^n(\bold{x}))d\bold{x}\bigg|\leq


\frac{R}{\sqrt{2}}\bigg\{\int\limits_


{B_R(\bold{y})}(\Rot\boldsymbol


{\varphi}_R^n(\bold{x}))^2d\bold{x}\bigg\}^{1/2}\bigg\{\int\limits_{B_R


(\bold{y})}\bigg(\int_0^1 (


\Rot\boldsymbol{\psi}_R^n(\bold{x}))d\tau\bigg)^2d\bold{x}\bigg\}^{1/2},


$$


откуда


$$


\bigg|\int_{\Omega}(\boldsymbol{\varphi}_R^n(\bold{x})\cdot\Rot


\boldsymbol{\psi}_R^n(\bold{x}))d\bold{x}\bigg|\leq\frac{R}{\sqrt{2}}\bigg\{\int_


{\Omega}(\Rot\boldsymbol


{\varphi}_R^n(\bold{x}))^2d\bold{x}\bigg\}^{1/2}\bigg


\{\int_{\Omega}


(\Rot\boldsymbol{\psi}_R^n(\bold{x}))^2d\bold{x}\bigg\}^{1/2}.


$$





Переходя в последнем неравенстве к пределу при $n\to\infty$ и принимая во


внимание \ref{un}, получим оценку


$$


\int_{\Omega}\bold{u}^2(\bold{x})d\bold{x}\leq\frac{\sqrt{3}


d(\Omega)}{4}\bigg\{\int_{\Omega}(\Rot\bold{u}(\bold{x}))^


2d\bold{x}\bigg\}^{1/2}\bigg\{\int _{\Omega}


\bold{u}^2(\bold{x})d\bold{x}\bigg\}^{1/2},


$$


из которой следует доказываемое неравенство.


\end{pf}








\begin{thebibliography}{99}





\bibitem{1}


Вейль Г. {\em Метод ортогональной проекции в теории потенциала}. В кн.:


Вейль~Г. Математика. Теоретическая физика. --


М.: Наука, 1984. -- 511~c.





\bibitem{2}


Ладыженская О.А.


{\em Математические вопросы динамики вязкой несжимаемой жидкости}. --


М.: Физматгиз, 1961. -- 203~c.





\bibitem{3}


Владимиров В.С.


{\em Обобщенные функции в математической физике}. --


М.: Наука, 1976. -- 280~c.





\bibitem{4}


Темам Р.


{\em Уравнение Навье--Стокса. Теория и численный анализ}. --


М.: Мир, 1981. -- 408~c.





\bibitem{5}


Дюво Г., Лионс Ж.-Л.


{\em Неравенства в механике и физике}. --


М.: Наука, 1980. -- 383~c.





\bibitem{6}


Лионс Ж.-Л., Мадженес Э.


{\em Неоднородные граничные задачи и их приложения}. --


М.: Мир, 1971. -- 371~c.





\bibitem{7}


Галанин М.П., Попов Ю.П. {\em Квазистационарные электромагнитные поля в


неоднородных средах}. --


М.: Физматгиз, 1995. -- 320~c.





\bibitem{8}


Быховский Э.Б., Смирнов Н.В. {\em Об ортогональном разложении пространства


вектор-функций, квадратично суммируемых по заданной


области, и операторах векторного анализа} // Тр. матем. ин-та АН


СССР. -- 1960. -- Т.\,59. -- С.\,5--36.





\bibitem{9}


Мазья В.Г.


{\em Пространства С.Л.\,Соболева}. --


Л.: Изд-во ЛГУ, 1985. -- 415~c.





\bibitem{10}


Калинин А.В.


{\em Некоторые оценки теории векторных полей} //


Вестн. Нижегородск. ун-та. Сер. матем. моделир. и


оптимал. управление. -- 1997. -- \No\,1. -- С.32--39.








\end{thebibliography}





\vskip 2cm





\post{\it Нижегородский государственный университет}{\it Поступила}


\post{}{10.01.2002}





\label{end}


\end{document}








